\section{Results}\label{sec:results}

\subsection{Historical Directional Evidence}\label{sec:historical-evidence}

Across all four measurement specifications, decreases in Wasserstein squared
distance are positively
associated with reference firms' relative weighted excess returns, and the central
article-redraw ranges remain positive
(Table~\ref{tab:directional-evidence}).
For a scale-free comparison within the primary specification, a difference of
one root-mean-square unit in centered weighted excess returns corresponds to
a fitted decrease in distance equal to \ppnum[3]{ir-normalized-alignment}
times the root-mean-square variation in centered distance changes.
This normalized alignment,
$\widehat\beta\operatorname{RMS}(\widetilde R)/\operatorname{RMS}(\widetilde A)$,
has square \ppnum[4]{ir-d2560-expanded-within-fit}, the in-sample fit of this
single-regressor, within-transformed projection.
It describes a small alignment in the observed sample.

For the primary specification, the coefficient is
\ppnum[6]{ir-primary-beta} squared-distance units per unit additive
weighted excess-return measure, with conditional article-redraw range
$[\ppnum[6]{ir-primary-low},\ppnum[6]{ir-primary-high}]$.
A \ppnum[2]{ir-financial-example} difference in the weighted excess-return measure
corresponds to a fitted approach
difference of \ppnume[2]{ir-primary-percent-effect} squared-distance units.
Because the coefficient describes direction in embedding geometry, its scale
differs across representations and has no percentage-capacity interpretation.

% AUTO-GENERATED from saved Paper 6 results; do not edit.
\begin{table}[htbp]
\centering
\caption{Complete-sample historical directional association. Native denotes 2,560 normalized coordinates. Percentiles describe conditional article-resampling sensitivity under the scheme in Section~\ref{sec:data}. Alignment is the normalized pooled association; its square is descriptive within-projection fit. All measurement specifications use 100 firms, fixed 2018 reference profiles and four transitions ending in 2019--2022.}
\label{tab:directional-evidence}
\begin{tabular}{lrrrrr}
\toprule
Specification & $\widehat\beta$ & \shortstack{2.5th\\percentile} & \shortstack{5th\\percentile} & \shortstack{97.5th\\percentile} & Alignment \\
\midrule
Native, expanded & \ppnum[6]{ir-d2560-expanded-beta} & \ppnum[6]{ir-d2560-expanded-low} & \ppnum[6]{ir-d2560-expanded-q05} & \ppnum[6]{ir-d2560-expanded-high} & \ppnum[4]{ir-d2560-expanded-alignment} \\
Native, matched & \ppnum[6]{ir-d2560-matched-beta} & \ppnum[6]{ir-d2560-matched-low} & \ppnum[6]{ir-d2560-matched-q05} & \ppnum[6]{ir-d2560-matched-high} & \ppnum[4]{ir-d2560-matched-alignment} \\
Leading 64, expanded & \ppnum[6]{ir-d64-expanded-beta} & \ppnum[6]{ir-d64-expanded-low} & \ppnum[6]{ir-d64-expanded-q05} & \ppnum[6]{ir-d64-expanded-high} & \ppnum[4]{ir-d64-expanded-alignment} \\
Leading 64, matched & \ppnum[6]{ir-d64-matched-beta} & \ppnum[6]{ir-d64-matched-low} & \ppnum[6]{ir-d64-matched-q05} & \ppnum[6]{ir-d64-matched-high} & \ppnum[4]{ir-d64-matched-alignment} \\
\bottomrule
\end{tabular}
\end{table}


All annual and deletion-based estimates remain positive.
The primary annual ratios are \ppnum[6]{ir-annual-2019-beta},
\ppnum[6]{ir-annual-2020-beta}, \ppnum[6]{ir-annual-2021-beta} and
\ppnum[6]{ir-annual-2022-beta} in 2019--2022, while leaving out one
transition gives coefficients from
\ppnum[6]{ir-year-deletion-low} to \ppnum[6]{ir-year-deletion-high}.

Removing each firm in all target and reference roles, then recentering over
remaining references, gives the primary range
$[\ppnum[6]{ir-role-deletion-low},\ppnum[6]{ir-role-deletion-high}]$.

These comparisons contribute unevenly to the pooled estimate:
the 2020 transition contributes \ppnum[1]{ir-primary-2020-weight-pct}\%
of the pooled second moment, and the largest reference share is
\ppnum[1]{ir-reference-share-pct}\%, compared with
\ppnum[2]{ir-target-share-pct}\% for a target.
A reference with unusually dispersed weighted returns can affect many target
comparisons, so equal reference weights need not imply equal influence on
the pooled slope.

The restricted movement profile improves held-out reconstruction of annual
news changes by only \ppnum[2]{ir-observation-primary-heldout-pct}\%
relative to predicting zero centered change, and two alternative
specifications perform worse than that benchmark.
This article-fold audit examines changes in news geometry separately from
their association with returns.
With the profile and reference ordering fixed from 2018 geometry, it fits an
annual movement magnitude on one article fold and evaluates it on the other.
Weak held-out reconstruction shows that annual article changes extend beyond
this local benchmark.
The measured movements also frequently reach or exceed the distance to the
nearest distinct reference distribution at the initial date.
These geometric departures do not identify crossings of operating positions,
which remain unobserved.
Appendix~\ref{app:observation-audit} reports the complete audit.
